% Generated draft: requires mathematical and source review.
\documentclass[11pt]{article}
\usepackage[margin=25mm]{geometry}
\usepackage{fontspec}
\setmainfont{Latin Modern Roman}
\usepackage{amsmath,amssymb,mathtools,unicode-math}
\setmathfont{Latin Modern Math}
\usepackage{enumitem,needspace,xurl,hyperref}
\setlength{\parindent}{0pt}
\setlength{\parskip}{5pt}
\setlength{\emergencystretch}{3em}
\newcommand{\N}{\mathbb{N}}
\newcommand{\Z}{\mathbb{Z}}
\newcommand{\Q}{\mathbb{Q}}
\newcommand{\R}{\mathbb{R}}
\newcommand{\C}{\mathbb{C}}
\newcommand{\F}{\mathbb{F}}
\newcommand{\A}{\mathbb{A}}
\newcommand{\PP}{\mathbb{P}}
\newcommand{\E}{\mathbb{E}}
\newcommand{\eps}{\varepsilon}
\newcommand{\dd}{\,\mathrm{d}}
\begin{document}
\section*{UnsolvedMath ID: 30005434}
\textbf{Source problem code:} catalogue entry
\subsection{Definitions and mathematical statement}
A semilattice is a set \(S\) with a binary operation \(\wedge\colon S\times S\to S\) satisfying, for all \(x,y,z\in S\),
\[
x\wedge y=y\wedge x,\qquad (x\wedge y)\wedge z=x\wedge(y\wedge z),\qquad x\wedge x=x.
\]
Its associated order is \(x\leq y\) if and only if \(x\wedge y=x\). An isomorphism of semilattices is a bijection preserving \(\wedge\). The supplied record does not specify whether empty semilattices are permitted or whether bounds or other additional conditions are required.

\textbf{Realization problem.} Characterize, up to semilattice isomorphism, the semilattices \((S,\wedge)\) for which there exists a \(\wedge\)-semibrace whose underlying semilattice is isomorphic to \((S,\wedge)\). Characterize also those for which such a realizing \(\wedge\)-semibrace can be chosen to be an \(L\)-algebra. Thus the requested characterization must give necessary and sufficient conditions for each of these two existence assertions.

\textbf{Unresolved definitions and source qualification.} The supplied record gives neither the axioms of a \(\wedge\)-semibrace nor those of an \(L\)-algebra, and does not specify precisely which operations or distinguished elements constitute their common structure. Consequently, the assertions above retain these source terms and are conditional on recovering their intended definitions; this draft does not supply a complete axiomatic formulation. The extracted text also remarks that there is no obvious notion of a trivial \(\wedge\)-semibrace.

The field labelled \texttt{original\_statement} does not contain the displayed realization question. Instead, it interleaves a passage labelled Problem 3 about self-similar closures with a corrupted definition of a right brick and questions labelled Problems 5 and 6 about underlying lattices of bricks. The closure passage attributes to Rump an embedding of each \(L\)-algebra \(X\) into a self-similar \(L\)-algebra \(S(X)\), generated as a monoid by \(X\), unique up to isomorphism, and possessing a group of left fractions
\[
G(X)=\{a^{-1}b:a,b\in S(X)\},
\]
called the structure group of \(X\). These are supplied contextual claims, not additional hypotheses of the realization problem.

The brick passage describes two \(L\)-algebra structures \((X,\cdot,e)\) and \((X,\circ,e)\), a further distinguished element \(0\), and additional identities whose extraction is corrupted. It asserts that the two canonical orders coincide and form a lattice. Its separate lattice questions concern \(L(K,n)\), the lattice of \(K\)-linear subspaces of the left vector space \(K^n\), where \(K\) is a possibly noncommutative field and \(n\geq1\), and the subspace lattice of a non-desarguesian plane. The extraction calls a lattice isomorphism to such an \(L(K,n)\) desarguesian and ends with an assertion that only positive answers are known, followed by an omitted list of cases. None of these separate questions or incomplete claims is substituted for the semilattice realization target.
\subsection{Short English statement}
Determine exactly which semilattices admit a realizing \(\wedge\)-semibrace, and which admit one that is also an \(L\)-algebra. The required algebraic axioms are absent, and the extracted original passage concerns different questions, so a fully precise formulation requires source reconciliation and human mathematical review.
\subsection{Sources}
\begin{itemize}
\item Local UnsolvedMath registry, supplied record with permanent ID \(30005434\), title \emph{Semilattice Realizations by Wedge Semibraces}, exact source code \texttt{catalogue entry}. The record gives proposed year 2023 and names no proposer. Its supplied status is \texttt{open}; its statement status is \texttt{corrected\_verified}. The accompanying claim that the official report was checked directly is supplied registry provenance, not verification performed for this draft.
\item \emph{Mini-Workshop: Skew Braces and the Yang--Baxter Equation} (2023), Oberwolfach Reports, pp.\ 537--563, DOI \texttt{10.4171/owr/2023/9}. Supplied source report and literature-source label: ``original OWR report.'' \url{https://doi.org/10.4171/owr/2023/9}. The registry describes Oberwolfach Reports as documenting mathematical workshops and their open-problem discussions. This authoritative report citation is supplied, but the relevant definitions and exact problem location are not reproduced reliably in the record.
\item Rump, \emph{L-algebras with duality and structure groups}. Supplied literature reference: \url{https://doi.org/10.1016/j.jpaa.2020.106314}. No additional bibliographic details are supplied.
\item Supplied literature assessment dated \textbf{2026-08-22}: the semilattice realization questions for \(\wedge\)-semibraces, including the \(L\)-algebra case, are labelled open; later work is described as providing canonical closures and finite enumerations without the requested characterization. This is supplied dated curation derived from the OWR record, not an independently verified literature assessment. The supplied background explicitly calls for independent verification of time-sensitive status and cited claims.
\end{itemize}
\end{document}
